\section{Dependent nested limits at every finite level}\label{sec:nested}
For $k\geq1$ define
\begin{equation}\label{eq:nested-definition}
 Q(M)=\lim_{L_1}\cdots\lim_{L_{k-1}}
       \Delta(M,L_1,\ldots,L_{k-1}),\qquad
 z_M=\frac{\log(1/Q(M))}{\log M}.
\end{equation}
The rightmost limit is taken first. Require every intermediate limit to
exist at each fixed tuple of remaining parameters. Require the final
$Q(M)$ to be finite and positive for every sufficiently large $M$;
earlier intermediate limits need not be positive because their
logarithms are not taken.

\begin{theorem}[Nested-box classification]\label{thm:nested}
For algebraic finite-range interactions, finite lower, upper and
existing exponents in \eqref{eq:nested-definition} have ceilings
$\mathsf S_{k+1},\mathsf P_{k+1},\mathsf D_{k+1}$ respectively.
Every ordered nonnegative finite pair in those endpoint classes is
realized by fixed integer onsite and nearest-neighbour interactions
in dimension $k+1$. The local dimension and the finite program may
depend on the pair and on $k$.
\end{theorem}
\begin{proof}[Necessity]
The finite spectra have uniform computable names. Applying the
$\mathsf D_m$ limit rule inductively gives a uniform
$0^{(k-1)}$-Cauchy name for $Q(M)$. Positivity permits the logarithm
in the same oracle. The final lower, upper or existing limit has the
asserted ceiling by Proposition~\ref{prop:calculus}.
\end{proof}

Here is a finite local-rule specification completing realization for
every fixed $k\geq3$. The cases $k=1,2$ are the fast finite-size
engine and Theorem~\ref{thm:film} respectively. Let $A$ be computable
in $0^{(k-1)}$. By the iterated limit lemma there is a total
computable binary function
\[
 A(j)=\lim_{s_1}\cdots\lim_{s_{k-1}}
 R(j,s_1,\ldots,s_{k-1}),
\]
with every intermediate bit limit existing. Set $R=1$ at $j=0,1,2$.
Use coordinates
\[
 (x,y,w_1,\ldots,w_{k-2},z),\qquad
 \text{side lengths }(L_{k-1},L_{k-1},L_1,\ldots,L_{k-2},M).
\]
The following fields
and predicates are finite for fixed $k$.

\begin{enumerate}
\item Give every axis two exact face flags, with the flag identity
of \eqref{eq:film-flags} in dimension $k+1$. For each
$a\in\{w_1,\ldots,w_{k-2},z\}$ introduce its own $(b_a,c_a,f_a)$
counter on every $(x,a)$ slice. Initialize $b_a=f_a=0$ at $a=0$,
set $c_a=1$ at $x=0$, and use the same carry and vertical recursion.
No copying along the other axes is needed: each slice is uniquely
determined from its own boundary initialization.
\item The only computation plane is the conjunction
$P=Z_0\bigwedge_i W_i^1$. On $P$ use the deterministic machine
compiler of Section~\ref{sec:classical}; off $P$ fix every tape,
head, output and arrival field to a unique dummy. Machine bonds
require $P$ at both endpoints and the two missing-neighbour arrival
conditions are imposed.
\item On $P$, expose the read-only pair of bits
$(b_{w_i},b_{w_i}\mathbin{\mathrm{OR}}f_{w_i})$ for each $i$.
At the far face the second track marks exactly the prefix of length
$\bitlength(L_i-1)$. Scanning to its first zero gives a delimiter,
and the first track then specifies $L_i-1$ in least-significant-first
binary order. Zero has its delimiter at the first cell. Merely
supplying an undelimited infinite string of zero-padded bits would
not let a finite-time machine recognize the input's end.
\item After parsing the inputs, the ordinary machine computes
$R(j,L_1-1,\ldots,L_{k-2}-1,s)$ for $j=0,\ldots,s$ at stage $s$
and writes the answers on its persistent output track. Every finite
computation finishes. A small box not exposing a delimiter still
has a unique finite computation prefix; completed output is only
asserted for sufficiently large innermost squares.
\item On $P_w=\bigwedge_iW_i^1$ transport the output from $Z_0$
along $z$, and fix the transport field to zero off $P_w$. The only
quantum segment is
$P_wY_1Z_1(b_z\mathbin{\mathrm{OR}}f_z)$. Its radix uses the
transported output; every other quantum site is frozen. Use the
same region guards, endpoint tests, penalties and weighted
transitions as the film. For eventual boxes the segment has an
$x$ neighbour outside the region, so wrap is still two-site.
\end{enumerate}

All tests are onsite or on an oriented nearest-neighbour bond.
Flags, independent counters, machine and transport are determined
in that order, which proves uniqueness of the correct classical
configuration. Any other classical block costs at least one;
its dynamic part is nonnegative. Thus the sector proof is unchanged
in higher dimension.

Fix $M$ and every outer size. For the finitely many read positions
$j<\ell(M)$, the innermost bit limit exists, and each bit is eventually
constant. The entire radix vector is therefore eventually constant
in $L_{k-1}$. Its gap is less than one once the first three bits are
read, and the innermost gap limit is precisely that full counter gap.
Let the next size grow and repeat: each next intermediate bit limit
again stabilizes a finite vector. Induction in the declared order gives
\[
 Q(M)=\Delta^+(K[2^{\beta A(j)}:j<\ell(M)]).
\]
This uses eventual constancy of finite discrete vectors, not continuity
of an eigenvalue through degeneracy or a uniform settling modulus.
The counter bracket gives
$z_M=2\beta\rho_{\ell(M)}(A)+o(1)$. Choose the target density pair
with Lemma~\ref{lem:relative-density} and an integer
$\beta\geq\max(1,b/2)$. This completes realization and the proof of
Theorem~\ref{thm:nested} at the finite-alphabet compiler level.

The dimension $k+1$ is a sufficient construction dimension. Minimality
is not proved. Nor does the theorem supply one local dimension for
all levels or all prescribed targets. A nondecreasing unbounded
sampling function that attains all sufficiently large integers
through unit steps preserves these density endpoints; arbitrary sparse
sampling need not do so.
